\documentclass[11pt]{article}
\usepackage[a4paper,margin=27mm]{geometry}
\usepackage{amsmath,amssymb,amsthm,mathtools,booktabs,array,hyperref,xurl,microtype}
\setlength{\emergencystretch}{3em}
\newtheorem{theorem}{Theorem}[section]
\newtheorem{lemma}[theorem]{Lemma}
\newtheorem{proposition}[theorem]{Proposition}
\newtheorem{corollary}[theorem]{Corollary}
\newtheorem{remark}[theorem]{Remark}
\newcommand{\Q}{\mathbb Q}\newcommand{\Z}{\mathbb Z}\newcommand{\N}{\mathbb N}
\newcommand{\A}{\mathcal A}\newcommand{\Tr}{\operatorname{Tr}}\newcommand{\Nm}{\operatorname{N}}\newcommand{\rad}{\operatorname{rad}}\newcommand{\ord}{\operatorname{ord}}
\newcommand{\OK}{\mathcal O_K}\newcommand{\vp}{v_p}\newcommand{\Sp}{\operatorname{Spec}}
\title{Intersective divisibility, local interpolation, and\protect\\valuation-tagged finite algebras\protect\\\large Extensions of Sun's ten-unknown construction for the $\Z\Q$ Diophantine Project}
\author{$\Z\Q$ Diophantine Project --- technical research note}
\date{24 November 2026}
\begin{document}
\maketitle
\begin{abstract}
Starting from Sun's arXiv:2609.01594v2, we prove an optimal degree-five
replacement for the degree-six intersective divisibility gadget, valid over
all number fields and rings of $S$-integers. The replacement allows exact
prescription of positive local valuations and unit parts. We give a uniform
valuation-isometric interpolation theorem for Lucas matrices and formal elliptic
groups. We sharpen the direct elliptic-index transfer in the project: for a
source of degree $D$, the quantitative modular-height argument needs only
precision $\lfloor D/2\rfloor+1$ in place of $3D+1$, with explicit constants
and a substantially shorter formal-logarithm polynomial. Finally, using
fixed prime-adic valuation tags, we give a direct simultaneous-square
polynomial and a guarded two-stage finite-algebra encoder. For the actual
21-point, two- or three-pole direct-index architectures, the latter has
single-branch total-degree bounds $2^{140}\delta$ and $2^{143}\delta$
(the full 11-branch bounds are multiplied by 11), where
$\delta$ is a precisely specified computable operator/coefficient-degree
majorant. It requires the same single encoder coordinate, no spectral
difference polynomial, and tolerates zero terminal roots and nonreduced
specializations. These are proved algebraic and conditional arithmetic
improvements, not new global variable-count records. All underlying elliptic, prime-pattern, modular-height and
pole-realization premises of the receiving project constructions are retained.
\end{abstract}

\tableofcontents

\section{Scope, attribution, and the receiving constructions}
Sun \cite{sun26} proves a ten-integer-variable universal coding theorem.
The project's preceding audit accepted the theorem subject to the minor
repairs recorded there. The construction uses a modified Lucas certificate,
residue-class interpolation, an intersective polynomial, and a single
homogeneous square/divisibility merger. The September source is distinct
from Sun's later prime-pattern papers. We use Theorem 1.1 of \cite{sun26}
as a coding theorem, not as an assertion that Lucas certificates over
$\Q$ have the same integer-index characterization. All new lemmas below
are independently proved, including both directions and exceptional fibers.
The ideas of intersectivity and merging, and the particular application to
Lucas interpolation, are due to Sun and the classical work he cites; no
priority claim is made for the classical intersective quintic itself.

The receiving repository \cite{repo} currently records $25/27/11$
coordinates for the unrestricted-degree integer definitions, and
$163/178/54$ for the quartics. The corresponding unrestricted-degree
undecidability frontier is $48/50/36$ existential coordinates, with the
first two using Sun's ten-integer source through the direct-index transfer
\cite{index}. The current quartic frontier is $884/1049/309$ coordinates
\cite{v124}. The last route in each list assumes ordinary rational $abc$;
all routes retain other arithmetic premises documented in their source
proofs. Counts for defining the integers include the input; undecidability
counts do not count the instance, which occurs in the coefficients.

This note has three distinct kinds of conclusions: unconditional
algebraic/number-theoretic lemmas; sharpened transfer theorems under the
explicitly retained premises of \cite{index}; and obstructions or
conditional opportunities for further coordinate savings. It does not
certify any new quartic table, remove the retained arithmetic hypotheses,
or identify a new sufficient numerical cutoff for the ineffective families.
The degree multipliers in Section~\ref{sec:receive} do not include a numerical
expansion of the (very large) coding polynomial or its input-dependent
coefficient DAG.

\section{An optimal intersective divisibility gadget}\label{sec:quintic}
Call $f\in\Z[X]$ \emph{intersective} if it has a root modulo every positive
integer. Equivalently it has a root in $\Z_p$ for every prime $p$, by CRT
and compactness. We require also that $f$ have no root in the ambient
field. Sun's Lemmas 5.1--5.2 use the monic sextic
$(X^2-13)(X^2-17)(X^2-221)$. The following quintic is sufficient.

\begin{theorem}\label{thm:quintic}
The polynomial
\[
 P(X)=(X^3-X-1)(X^2-X+6)=X^5-X^4+5X^3-5X-6
\]
has no rational root and has a simple root modulo every rational prime.
No nonconstant polynomial of degree less than five with rational
coefficients can have a root in every $\Q_p$ but no root in $\Q$.
\end{theorem}
\begin{proof}
The cubic is irreducible by the rational-root test (its only possible
integer roots are $\pm1$), and its discriminant is $-23$. The quadratic
has the same discriminant and is irreducible. At $p=2$ the quadratic has
simple roots $0,1$ and the cubic has neither. At $p=23$ the cubic has the
simple root $3$ (its derivative there is $26$), at which the quadratic has
value $12$ modulo 23. At any other odd prime the cubic is separable. If
it has no linear factor, its Frobenius permutation is a 3-cycle, which
is even; hence the discriminant is a nonzero square modulo $p$, and the
quadratic splits. If the cubic has precisely one linear factor, the
permutation is a transposition, the discriminant is a nonsquare, and that
linear root is not a quadratic root. If the cubic splits, it has three
distinct roots and at most two of them can also be quadratic roots. In
each case at least one root of the product is simple. The discriminant
criterion follows from the alternating Vandermonde product of the three
roots: Frobenius fixes it precisely for an even permutation. Hensel's
lemma now gives a $\Z_p$ root at every prime.

For minimality, work more generally over a number field $K$. If a
polynomial with no $K$-root has degree at most four, its distinct
irreducible factors are either a single irreducible factor of degree
2, 3 or 4, or at most two quadratics. An irreducible factor of degree
at least two has a transitive Galois action on its roots; the average
number of fixed points is one, while the identity fixes more than one.
Thus some group element has no fixed point. For two quadratic factors,
the compositum of their splitting fields has an automorphism nontrivial
on both quadratic fields (also when the fields coincide). This gives
an element with no fixed root of either factor. Chebotarev's density
theorem supplies an unramified prime of this Frobenius conjugacy class,
outside the finite set of coefficient denominators, leading coefficients
and discriminants of the distinct irreducible factors. No root exists
in the corresponding completion: after making the factors monic with
integral coefficients, any local root would be integral and reduce to
a fixed root of Frobenius. Repeated factors do not change this argument.
This proves the lower bound, over $\Q$ and indeed over every number field.
\end{proof}

\begin{theorem}[Number-field version]\label{thm:nf}
For every number field $K$ there is an effectively constructed monic
quintic $P_K\in\Z[X]$ with no root in $K$ and a simple root modulo every
rational prime. In particular it is intersective in $\OK$ and has a root
in every nonarchimedean completion of $K$.
\end{theorem}
\begin{proof}
Choose a rational prime $\ell\equiv1\pmod {12}$ unramified in $K$;
Dirichlet's theorem gives such a prime, and trial search is effective.
The integer 3 is a nonzero square modulo $\ell$ by quadratic reciprocity.
Consequently $4-27a^2\equiv0\pmod\ell$ has a solution $a$ not divisible
by $\ell$. Lift the solution modulo $\ell^2$ so that
$v_\ell(4-27a^2)=1$, possible since the derivative $-54a$ is a unit.
By CRT take $a$ odd, and put
\[
 f_a=X^3-X-a,\qquad D=4-27a^2,\qquad q_a=X^2-X+(1-D)/4.
\]
The integer $(1-D)/4$ is well defined. The cubic is irreducible over
$\Q$ by reduction modulo 2; its discriminant is $D$. The field generated
by a cubic root has discriminant divisible by $\ell$ (the index square
in the polynomial-discriminant identity cannot remove a valuation of
one), hence is ramified at $\ell$. The discriminant quadratic field
$\Q(\sqrt D)$ is also ramified at $\ell$. Neither field can be contained
in $K$, since ramification indices multiply in towers. Thus neither
factor has a root in $K$.

At an odd prime $p\mid D$, the cubic has a double root and a distinct
simple root in $\mathbb F_p$. Indeed $p\ne3$, and a triple root would
force its $X$ coefficient to vanish; the unique repeated root is fixed
by Frobenius, and so is the remaining root. For each such prime choose
$c$ modulo $p$ so that this simple root is not a root of $q_a(X-c)$.
There is only one forbidden class because $q_a$ has a double root modulo
$p$. Choose one integer $c$ satisfying these finitely many conditions.
Set $P_K=f_a(X)q_a(X-c)$. At odd primes not dividing $D$ the three
Frobenius cases in Theorem~\ref{thm:quintic} still apply, independent of
the translation $c$. At 2 the cubic is irreducible; since $a$ is odd,
$D\equiv1\pmod8$, and the quadratic has two simple roots modulo 2.
At primes dividing $D$ the chosen translation preserves a simple root
of the product. This proves the assertion. A $\Z_p$ root maps to a root
in each $\mathcal O_{K,\mathfrak p}$ over $p$, so CRT in $\OK$ supplies
a root modulo every nonzero ideal.
\end{proof}

\begin{lemma}[Exact local preparation]\label{lem:exact}
Let $P$ be either quintic above. For any finite set of rational primes
$p$, integers $e_p\ge1$, units $u_p\in\Z_p^\times$, and precision
$k_p\ge1$, there are infinitely many $y\in\Z$ with
\[
 P(y)\equiv p^{e_p}u_p\pmod {p^{e_p+k_p}}\quad\text{for every selected }p.
\]
One can require $y$ arbitrarily large of either sign. A congruence for
$y$ at other primes may be appended. In a number field one may similarly
prescribe values in $\mathfrak p\mathcal O_{K,\mathfrak p}$ to finite
precision at distinct prime ideals, with $y\in\OK$.
\end{lemma}
\begin{proof}
At a simple residue root $r$, the polynomial $P(X)-t$ for any
$t\in p\Z_p$ has a unique $\Z_p$ root congruent to $r$ modulo $p$.
Use $t=p^{e_p}u_p$, retain its residue to the prescribed precision,
and apply CRT. The resulting integer progression is unbounded in
both directions. The proof over $K$ is identical in each local ring,
since the derivative at the lifted simple rational residue root
remains a unit; use the ideal CRT. Positive depths are essential:
for example $P(y)$ is even for every integer $y$, so arbitrary depth
zero at 2 cannot be claimed. No restriction at unselected primes is
implied.
\end{proof}

\begin{theorem}[Divisibility with a nonzero guard]\label{thm:div}
Let $R=\mathcal O_{K,S}$ for a finite set of inverted finite primes,
and use $P_K$ from Theorem~\ref{thm:nf} (use $P$ over $\Q$). Define the
homogenization $F(X,S_0)=S_0^5P_K(X/S_0)\in\Z[X,S_0]$. For
$S_0,T,\xi\in R$, $S_0\ne0$, one has
\begin{equation}\label{eq:div}
 (S_0\mid_R T\text{ and }\xi\ne0)
 \ \Longleftrightarrow\ \exists Y\in R:\quad
 \xi S_0^5\mid_R F(T+S_0Y,S_0).
\end{equation}
More generally, for $x\in K$ and $\xi\in R$,
\begin{equation}\label{eq:rational}
 x\in R,\ \xi\ne0
 \quad\Longleftrightarrow\quad
 \exists Y,Z\in R:\ P_K(x+Y)=\xi Z.
\end{equation}
\end{theorem}
\begin{proof}
For $S_0\ne0$ the homogeneous value never vanishes on $K$ because
$P_K$ is rootless. Thus the divisibility on the right in \eqref{eq:div}
excludes $\xi=0$, under the convention $a\mid b$ iff $b=az$ for
some $z\in R$. If $v_\mathfrak p(T)<v_\mathfrak p(S_0)$ at an uninverted
prime then $v_\mathfrak p(T+S_0Y)=v_\mathfrak p(T)$. Since $P_K$ is monic
with integral coefficients, $v_\mathfrak p(P_K(w))=5v_\mathfrak p(w)$
when $v_\mathfrak p(w)<0$. Therefore $F$ has valuation
$5v_\mathfrak p(T)$, less than $v_\mathfrak p(\xi S_0^5)$. This
contradicts divisibility. Conversely, if $x=T/S_0\in R$ and
$\xi\ne0$, select a root of $P_K$ modulo $\xi R$ by CRT in $R$.
Translate it by $-x$ to obtain $Y$ and multiply by $S_0^5$.
The same monicity argument, now applied directly to $x+Y$, proves
\eqref{eq:rational}. Conversely choose the same residue root and set
$Z=P_K(x+Y)/\xi\in R$.
\end{proof}

The formula remains valid if an existing integral witness $Y$ is free
to range through a progression $Y_0+MR$ with $(M,\xi)=R$; CRT chooses
the required residue of $Y$ modulo $\xi$. More generally one only needs
the intersection of the two residue conditions to be nonempty. If several
independent shifts are constrained to be the same, pairwise coprime
moduli suffice; overlapping moduli need an actual compatibility check.
Thus the replacement is degree-optimal within the class of monic rootless
intersective univariate polynomials. Whenever the receiving use of Sun's Lemma 5.2 requires only its
equivalence and his Theorem 4.1 permits an arbitrary prescribed residue
class, the same residue-selection argument works for the quintic.
The replacement pair is $(\xi S_0^5,F(T+S_0Y,S_0))$ in place of
$(\xi S_0^6,R_{\rm old})$. This saves one in the degree of the
homogeneous gadget without adding a witness. The full Section 7
coefficient source has not been regenerated; no new global degree
record or minimum for the final universal polynomial is claimed.

\begin{remark}[Sharp limitations]
If $Y$ is allowed to range over $K$ without the integrality requirement,
\eqref{eq:rational} fails: over $\Q$, take $x=1/2$, $Y=-1/2$, $\xi=1$,
$Z=P(0)=-6$. If $\xi$ is not integral at the relevant prime, take
$x=1/2$, $Y=0$, $\xi=P(1/2)$, $Z=1$. If $S_0=0$, the homogeneous
formula admits the spurious branch $T=0$ regardless of $\xi$.
The missing $S_0\ne0$ hypothesis must therefore be supplied by the
receiving circuit. Further, $P(X)=0$ has points in every finite
completion and a real point (the cubic has odd degree), but no rational
point. It is a particularly economical counterexample to replacing
global rational realizability by finite local tests. This is relevant to
the project's physical-lift and pole-realization problems; it does not
replace the Bhargava--Burungale global point theorem.
\end{remark}

\paragraph{Quartic accounting.}
The new quintic need not save quadratic registers. For example a direct
quadratic compilation of $F(T+S_0Y,S_0)=\xi S_0^5 Z$ uses
$X=T+S_0Y$, $s_2=S_0^2$, $b=X^2-s_2$, $c=Xb-S_0s_2$,
$d=X^2-XS_0+6s_2$, $s_4=s_2^2$, $s_5=S_0s_4$, $w=\xi Z$,
and $cd=s_5w$. These are eight private registers beyond $Z$.
The original sextic admits seven: $X,s_2,x_2=X^2$,
$b=(x_2-13s_2)(x_2-17s_2)$, $s_4=s_2^2$, $s_6=s_2s_4$,
$w=\xi Z$, and $b(x_2-221s_2)=s_6w$. The degree improvement is
therefore not a quartic-variable saving by this straightforward compiler.

\section{Valuation-isometric interpolation}\label{sec:interp}
Sun's Section 4 obtains full residue-class interpolation by choosing a
Lucas recurrence at suitable indices. The following local formulation
isolates the algebraic mechanism, applies to other recurrences, and
makes explicit the compatibility requirements for elliptic applications.

\begin{theorem}[Matrix and algebra version]\label{thm:matrix}
Let $B$ be a finite free $\Z_p$-algebra (not necessarily commutative),
$g=1+p^d a\in B$, $d\ge1$ for odd $p$ and $d\ge2$ for $p=2$.
Let $L:B\to\Z_p$ be linear with $L(a)$ a unit. Then
\[
 f(n)=p^{-d}L(g^n-1),\qquad n\in\Z,
\]
satisfies $v_p(f(n)-f(m))=v_p(n-m)$ for $n\ne m$. It induces a
permutation of every $\Z/p^k\Z$, and extends to a bijective isometry
of $\Z_p$. The result also holds with $L(x)$ replaced by $L(bx)$
provided $L(ba)$ is a unit.
\end{theorem}
\begin{proof}
The element $g$ is a unit by the convergent geometric series. For $h>0$,
\[
 g^h-1=h p^d a+\sum_{j=2}^h\binom hj p^{dj}a^j.
\]
The elementary identity $j\binom hj=h\binom{h-1}{j-1}$ gives
$v_p(\binom hj)\ge v_p(h)-v_p(j)$. For $j\ge2$,
$(j-1)d-v_p(j)\ge1$ in the stated cases. Indeed for odd $p$,
$p^{j-1}>j$ for $j\ge2$; for $p=2$ the factor $2(j-1)$ is sufficient.
The remainder lies in $p^{d+v_p(h)+1}B$. For $n>m$, multiply by $g^m$;
its reduction modulo $p$ is one even when $m$ is negative. Applying
$L$ yields the asserted exact valuation. The other order follows by
negation. The isometry descends to an injection and hence a bijection
on each finite residue ring; inverse limits give the extension.
\end{proof}

For the Lucas sequence $u_0=0,u_1=1,u_{n+1}=A u_n-u_{n-1}$ put
$M_A=\left(\begin{smallmatrix}A&-1\\1&0\end{smallmatrix}\right)$.
Its lower-left entry in $M_A^n$ is $u_n$ for every $n\in\Z$.
If $M_A^r=I+p^d B_0$ and the lower-left entry of $M_A^b B_0$ is a
unit, Theorem~\ref{thm:matrix} shows that
\[
 n\longmapsto \frac{u_{b+rn}(A)-u_b(A)}{p^d}
\]
is a permutation modulo every $p^k$. This gives simultaneous
interpolation across distinct primes by CRT, even with different
recurrences and starting indices (the common parameter $n$ is used in
each recurrence). The unit condition is indispensable; the theorem
is not a replacement for Sun's proof that his actual parameters satisfy
it. The depth hypothesis at 2 is sharp in general: $g=-1$, $d=1$,
$a=-1$ gives a two-valued function.

\begin{theorem}[Elliptic version]\label{thm:ellinterp}
Let $E$ have good reduction over a finite extension $L/\Q_p$, with
valuation normalized by $v(p)=1$. Suppose $Q$ is a nonzero point in
the formal subgroup with $d=v(z(Q))>1/(p-1)$, where $z=-x/y$ is a
minimal integral formal parameter. Let $R$ be any point in the formal
subgroup (including the identity). Then for distinct integers $n,m$,
\[
 v\bigl(z(R+[n]Q)-z(R+[m]Q)\bigr)=d+v_p(n-m).
\]
For $L=\Q_p$, the normalized map
$(z(R+[n]Q)-z(R))/p^d$ permutes every residue ring $\Z/p^k\Z$.
For general $L$ its image is a single $\Z_p$-parameter family, not
necessarily all of $\mathcal O_L$ after normalization.
\end{theorem}
\begin{proof}
The formal logarithm is an isometric group isomorphism on points of
parameter depth greater than $1/(p-1)$: its invariant differential
has integral coefficients and the term of degree $j\ge2$ has valuation
at least $jd-v_p(j)>d$. The formal inverse exists on this depth
ball, by the usual inverse power series or the strict contraction
of the higher terms. Thus $v(z([h]Q))=d+v_p(h)$ for nonzero $h$.
The formal group law satisfies $F(X,Y)-X=Y(1+H(X,Y))$ with
$H\in (X,Y)\mathcal O_L[[X,Y]]$ because its linear part is $X+Y$.
Take $X=z(R+[m]Q)$ and $Y=z([n-m]Q)$ in the maximal ideal to
obtain the displayed valuation. The permutation proof is as in
Theorem~\ref{thm:matrix}. Over an unramified extension of residue
degree $f>1$, the residues of a single integer parameter modulo $p$
have at most $p$ values, not $p^f$, proving the limitation.
\end{proof}

For a global rank-one point $P$, at finitely many good places let
$r_\mathfrak p$ be the order of its reduction and choose a common
multiple $r$ such that $[r]P$ has the required depth at every place.
Then the normalized scalar-multiple maps above interpolate arbitrary
integer-parameter residues at distinct underlying rational primes.
At two places over the same rational prime the congruences on the
single multiplier must agree modulo the smaller modulus. If the
reference pair is $[h]P,[h-1]P$, high depth at both split places
is not automatic: when the reduction orders are prime to $p$ and
$z([r_\mathfrak p]P)$ has depth one in both slots, depth at least two
would force $p\mid h$ and $p\mid h-1$ respectively, which is impossible.
This explains why replacing logarithm precision by deeper reference
points would require new arithmetic premises rather than a CRT trick.

\section{Shorter logarithms and sharper source precision}\label{sec:precision}
We now use the direct-index source \cite{index} as a fixed receiving
construction. We give details because changing the precision in this
source affects the height contradiction and completeness simultaneously.
Let $Q(a,Z_1,\ldots,Z_k)$ be the fixed integer coding polynomial;
allow also a finite system $Q_1=\cdots=Q_s=0$. The source uses a common
shift $\sigma\in\{0,\ldots,k\}$ so that $m_i=Z_i+\sigma$ are nonzero,
$m_0=a+1$, and
\[
 \Psi_{\sigma,j}(M_0,\ldots,M_k)=Q_j(M_0-1,M_1-\sigma,\ldots,M_k-\sigma).
\]
Set $D_j=\max(1,\max_\sigma\deg\Psi_{\sigma,j})$ and
$C_j=\max(1,\max_\sigma\|\Psi_{\sigma,j}\|_1)$, where the latter is
the sum of absolute integer coefficients. Write
$B=\max(1,|m_0|,\ldots,|m_k|)$. If a residual $N_j=\Psi_{\sigma,j}(m)$
is nonzero then $1\le|N_j|\le C_j B^{D_j}$.

\subsection{Logarithm truncation}
On the project's curve over $K=\Q(\sqrt2)$, with $D_E=11-6\sqrt2$,
the formal logarithm of $z=-x/y$ is \cite{index,ellheight}
\[
 \ell(z)=\sum_{r\ge0}\frac{(-1)^r\binom{2r}{r}D_E^r}{4r+1}z^{4r+1}.
\]
One may derive it by putting $w=1/x$, so $z^2=w/(1-D_Ew^2)$
and differentiating; the invariant differential is
$dz/(\sqrt{1+4D_Ez^4})$, whose binomial expansion gives the series.
For clarity, differentiating $z=-x/y$ and using $y^2=x^3-D_Ex$
gives $dz/(dx/(2y))=(x^2+D_E)/(x^2-D_E)$, whose square is
$1+4D_Ez^4$ and whose value at the identity is one.

\begin{lemma}[Short truncation]\label{lem:short}
Let $L\ge1$ and $p_0\ge3$ be an integer lower bound for the odd
unramified good residue primes under consideration. Take
\begin{equation}\label{eq:cutoff}
 r_0=\min\{r\ge1: 4r-\lfloor\log_{p_0}(4r+1)\rfloor\ge L\},\qquad
 \ell^{[L]}(z)=\sum_{r=0}^{r_0-1}
 \frac{(-1)^r\binom{2r}{r}D_E^r}{4r+1}z^{4r+1}.
\end{equation}
For $v_p(z)\ge1$ and $p\ge p_0$, $z\ne0$, the exact ratio
$\ell^{[L]}(z)/\ell(z)$ belongs to $1+p^L\mathcal O_v$.
The cutoff is computable by integer powers alone. One has
$r_0\le\lceil L/3\rceil$; if $p_0>4\lceil L/4\rceil+1$, then
$r_0=\lceil L/4\rceil$. The polynomial degree is at most $4r_0-3$.
\end{lemma}
\begin{proof}
The relative term of index $r\ge1$ has valuation at least
$4r-v_p(4r+1)\ge4r-\lfloor\log_{p_0}(4r+1)\rfloor$.
The integer function $n-\lfloor\log_{p_0}(n+1)\rfloor$ is
nondecreasing: at each step the floor increases by at most one.
Also $3^{r+1}>4r+1$ for $r\ge1$ by induction, so the bound is
at least $3r$. Thus $\ell(z)/z$ is a unit and the omitted relative
terms have valuation at least $L$. This proves the ratio estimate
and the first cutoff bound. If the stated strict inequality for
$p_0$ holds, the floor is zero at $r=\lceil L/4\rceil$; smaller $r$
cannot work because $4r<L$.
\end{proof}

\subsection{The direct-index residue test}
The point charts, guards and coefficients are those of \cite{index}:
$A_i=[m_i]P$, $U_i=[j_i]P$, common $H=[h]P$, with finite nonzero
points and $h\ne0,\pm1$. Write $f(Q)=1/x(Q)$, $z(Q)=-x(Q)/y(Q)$,
and for a contact argument $t$ put
\begin{align*}
 \Theta&=\{0,-1/2,-1,1/2399,-32/407\},\qquad
 \Delta=t+1,\\
 G(t)&=\prod_{\theta\in(2+\Theta)\cup(4+\Theta)}(t-\theta),\qquad u=G(t)/\Delta^{10}.
\end{align*}
Include the existing exclusions (bad primes, coefficient primes,
contact differences and $\Delta$ at contact, ramified primes, etc.)
in the fixed finite set $S$. All divisions and nonzero guards of the
source remain; zero comparison and terminal test values are permitted.
Take precision $L=\max_j L_j$. Use
\[
 r_i=\frac{\ell^{[L]}(z(U_i))}{\ell^{[L]}(z(H))},\qquad
 r'_i=\overline{\frac{\ell^{[L]}(z(U_i-A_i))}{\ell^{[L]}(z(H-P))}}.
\]
The input list of field-valued tests $C$ comprises, each divided by
$\Delta^2$, the following entries:
\begin{align*}
 &f(H),\ f(U_i),\ \overline{f(H-P)},\ \overline{f(U_i-A_i)},\\
 &(r_i-r'_i)/u^{L-1},\qquad
 \Psi_{\sigma,j}(a+1,r_1,\ldots,r_k)/u^{L_j-1}.
\end{align*}
Duplicate tests can be omitted. The project uses the exact product gate
\begin{equation}\label{eq:gate}
 D_C=\frac{C^2}{1+uC^2},\quad W=\prod_C(1+D_C^2/u),\quad V=\Tr(W^{-1}),\quad
 e=7u\left(\frac{16}{101(1+u)V}\right)^2,
\end{equation}
with terminal normalized BB product pole
$q^2=(1-e)(v^4-(7u)^2)(w^4-(7u)^2)$.
Its unconditional local soundness is included here to specify exactly
what the changed tests require. At an unramified odd consumed prime
$p\notin S$ with $\nu=v_p(u)>0$ odd, both competing denominator
valuations in $D_C$ differ in parity, so $v_p(D_C)=\max(2v_p(C),-\nu)$.
The same unequal-parity argument gives $v_p(1+D_C^2/u)=\min(0,2v_p(D_C)-\nu)$.
The product pole forces $e$ not to have negative odd valuation:
otherwise $v_p(7u)$ is odd, each quartic factor has even valuation
without cancellation, and $1-e$ has odd valuation. Since $v_p(e)$
is odd at the consumed prime, it is strictly positive. Writing $q_0$
for the maximum of the slot valuations of $W$, $v_p(V)\ge-q_0$, and
\[
 0<v_p(e)\le \nu+2q_0.
\]
In a slot attaining $q_0$ every factor valuation is $>-\nu/2$, hence
$v_p(D_C)>\nu/4$ and $v_p(C)>0$ for every test. The reciprocal tests
place the required points in their formal subgroups (in the relevant
split slots or at the unramified place). At such a contact $\Delta$
is a unit. The formal logarithm is additive, and by multiplying the
base point into the formal subgroup its exact ratios are $j_i/h$ and
$(j_i-m_i)/(h-1)$. Lemma~\ref{lem:short} gives
\[
 h r_i\equiv j_i,\qquad (h-1)r'_i\equiv j_i-m_i\pmod {p^L}.
\]
This assertion does not require $r_i$ to be integral initially: the
relative errors are multiplied by the integer indices. The comparison
imposes $r_i-r'_i\equiv0\pmod {p^L}$, because $(L-1)\nu+1\ge L$.
Subtracting gives $r_i\equiv m_i\pmod {p^L}$, including integrality
of $r_i$ in that slot. Evaluating the integer polynomial then yields
$p^{L_j}\mid N_j$ for every $j$. The unramified hypothesis allows
this conclusion over $\Z$. The product $M_S(t)$ of all the consumed
positive odd primes therefore satisfies
\begin{equation}\label{eq:raddiv}
 M_S(t)^{L_j}\mid N_j\quad\text{for each }j.
\end{equation}

\subsection{A sharper height contradiction}
The retained quantitative modular input in \cite{index} is
\begin{equation}\label{eq:modular}
 h(t)\le 2^{551} M(t)(199+\tfrac12\log M(t))^2,
\end{equation}
where $M(t)$ is its consumed-prime radical with the theorem's own
exclusions. Enlarge $S$ to include those exclusions as in the original
transfer. If $R_S$ is the product of primes in $S$ (or an integer upper
bound), $M(t)\le R_S M_S(t)$. This is a retained premise, not an
independent proof of the modular estimate.

\begin{theorem}[Improved effective precision]\label{thm:precision}
For each source residual of degree bound $D_j$ one may take
$L_j=\lfloor D_j/2\rfloor+1$. Put $b(N)=\lfloor\log_2 N\rfloor+1$
for positive integers, and $C'_j=\lceil C_j^{1/L_j}\rceil$.
If $N_j\ne0$, the tests imply $h(t)\le H_j B^2$, where the explicit
integer constant
\begin{equation}\label{eq:H}
 H_j=2^{551} R_S C'_j(199+b(R_S)+b(C'_j)+2L_j)^2
\end{equation}
is independent of the instance $a$ and the witnesses. The same
48/50-variable direct-index transfers follow with the shorter
truncation of Lemma~\ref{lem:short}, under their retained arithmetic
premises; the three-contact sufficient exponent is effective.
\end{theorem}
\begin{proof}
By \eqref{eq:raddiv}, $M_S(t)\le C'_j B^{\alpha}$ with
$\alpha=D_j/L_j<2$ and $2-\alpha=\eta\ge1/L_j$.
Use this bound in \eqref{eq:modular}; with
$A=199+b(R_S)+b(C'_j)$ the result is at most
$2^{551}R_SC'_j B^\alpha(A+\log B)^2$.
For $x=\log B\ge0$,
$e^{-\eta x/2}(A+x)\le A+2/\eta\le A+2L_j$ (maximize
$x e^{-\eta x/2}$, or use the elementary exponential bound).
Squaring proves \eqref{eq:H}.

Let $H_* = \max_j H_j$. If some source residual is nonzero, the bound
holds at every contact argument with this common constant. For the
two-contact transfer use $t_0=\lambda$, $t_1=\lambda+T^J$ with
$T=\Tr(\sum_{i=0}^k x(A_i))>1$. The source's aggregate trace-height
premise gives $B^2\le A_*h(T)$ with a fixed $A_*$ after absorbing its
additive constant (as in \cite{index}). Since
$J h(T)\le h(t_0)+h(t_1)+\log2\le 2H_*B^2+\log2$, any fixed integer
$J>A_*(2H_*+1)$ contradicts $B^2\le A_* h(T)$.
For the specified three-contact route write
$Z=\sum_{i=0}^k x(A_i)=U+V_1\sqrt2$ and take
$t_0=\lambda,t_1=\lambda+U^J,t_2=\lambda+V_1^J$.
The source supplies $B^2\le K_r(h(U)+h(V_1))$ with
$r=k+1$, $a_r=(r-1)2^{r-1}$ and $K_r=18+a_r^2+2r2^{r-1}$.
Hence $J B^2\le K_r(4H_* B^2+2\log2)$, contradicted by the effective
even integer $J=2K_r(4H_*+3)$. It follows that every residual is zero.

For completeness fix an integer solution and the common shift,
choose $A_i$ and the aggregate before selecting primes, and set
$U_i=[hm_i]P$ as in the source. At selected primes, all successful-slot
reciprocal tests are positive and the exact log ratios are $m_i$;
Lemma~\ref{lem:short} makes the comparison positive after division
by $u^{L-1}$ and the polynomial tests positive after division by
$u^{L_j-1}$ (the selected contact valuations equal one).
The source's fixed-place/real approximation first chooses $H$ away
from the finitely many zeros and poles of the new nonzero rational
logarithm functions, then preserves their bounded neighborhoods.
A nonconstant point map composed with a nonzero rational function
has only finitely many zeros on the elliptic curve; the reference
logarithms are nonzero rational functions because their linear term
is $z$. Include their coefficient denominator primes in $S$; the
shorter polynomial uses a prefix of the old coefficient list, so
this introduces no additional primes beyond those already excluded
in the previous precision-$3D+1$ transfer. At all other places the
source's parity and preparation calculations for \eqref{eq:gate}
are unchanged: they use boundedness at fixed places, the selected
unit trace, and the square factor, not the value of $L$. The BB
prepared pole theorem supplies the terminal solutions. This is
exactly the source's order of witness choices, with the same
ordinary coordinates and fewer logarithm terms. All its prime-pattern,
recurrence, modular and height premises are still required.
\end{proof}

If only an ineffective power estimate $h(t)\le H M_S(t)^E$ is retained,
with fixed $E>0$, use any integer $L_j\ge1$ satisfying $L_j\ge ED_j/2$:
then $h(t)\le H (C'_j)^E B^2$. This improves the old sufficient
precision $3D+1$ when its exponent was $E=3$. In the effective
logarithmic estimate \eqref{eq:modular}, strict inequality $D_j/L_j<2$
is essential to this particular argument: at equality the square of
$\log B$ is not absorbed by a fixed $B^2$ constant. For a source system,
testing each residual individually rather than a single sum of squares
reduces the precision determined by its degree. Different $L_j$ also
reduce the powers of $u$ in the individual tests. The tradeoff is the
number of factors of $W$; a global determinant-degree comparison needs
its coefficient DAG, not only the precision bound.

\section{Valuation-tagged simultaneous-square polynomials}\label{sec:square}
Sun's Lemma 2.6 uses a Matiyasevich--Robinson polynomial to merge square
conditions and divisibility. Here a fixed prime-adic tag replaces the
large variable-dependent positional base. This also works for rational
radicands, including negative values and zero roots.

\begin{theorem}[Tagged radical merger]\label{thm:square}
For $k\ge1$ choose an odd integer $m\ge k$ and a rational prime $p$.
Define
\begin{equation}\label{eq:J}
 J_{k,m}(A_1,\ldots,A_k,Z)=
 \prod_{\varepsilon\in\{\pm1\}^k}
 \left(Z-\sum_{i=0}^{k-1}\varepsilon_{i+1}p^i
 A_{i+1}^{(m-1)/2}\sqrt{A_{i+1}}\right).
\end{equation}
It belongs to $\Z[A_1,\ldots,A_k,Z]$, is monic of degree $2^k$ in $Z$,
and has total degree at most $2^k\max(1,m/2)$. For every rational tuple
$A$, it has a rational root if and only if all $A_i$ are rational squares.
For integral $A$ the equivalent root can be taken integral, and any
rational root is integral. Moreover the signed sum separates the
\emph{distinct} geometric characters of the tensor square-root algebra,
even in specializations with zero radicands.
\end{theorem}
\begin{proof}
Changing any formal radical's sign permutes the factors, so in the
expansion every radical exponent is even. This proves polynomiality,
monicity and the degree bound, using weight $m/2$ per radical term.
Fix a rational specialization, and choose one extension of the
$p$-adic valuation to the splitting field. The difference of the signed
sums at two distinct characters is a nonempty sum of terms
$\pm2 p^i A_{i+1}^{(m-1)/2}\sqrt{A_{i+1}}$ with nonzero radicands.
Their valuations are
\[
 v_p(2)+i+(m/2)v_p(A_{i+1}).
\]
Equality for different indices would give $2(i-j)\equiv0\pmod m$.
Since $m$ is odd and $|i-j|<m$, this is impossible. The unique term of
least valuation prevents cancellation. Thus the signed sum separates
characters. If its value is rational, every absolute Galois conjugate
has the same value, so separation fixes the character and all its
radicals are rational. Conversely rational radicals plainly give a
rational root. Rational roots of monic integral quadratics are integral,
which proves the integer assertion. Sign changes on zero roots do not
change the geometric character, and the proof never assumes reducedness
of the tensor algebra.
\end{proof}

The same proof holds over any number field using a fixed prime ideal
with normalized discrete valuation and an element of valuation one
in place of $p$, with coefficients in $\OK$ if this element is integral.
The Galois action is then over $K$. Anisotropic coefficient aggregation
in later applications over nonreal $K$ requires the modification in
Section~\ref{sec:guard}, not an unqualified sum of squares.

\begin{corollary}[Homogeneous merger]\label{cor:H}
For integers $S,T,M$ and $A_i$, $S\ne0$, put
\[
 H_{k,m}(A,S,T,M)=S^{2^k}J_{k,m}(A,M-T/S).
\]
Then $\exists M\in\Z:H_{k,m}=0$ is equivalent to the conjunction
$A_i\in\Z^2$ for all $i$ and $S\mid T$. The formula extends polynomially
to $S=0$ with value $(-T)^{2^k}$. For $k=4$, $m=5$, one may use
$\deg J_{4,5}\le40$ and $\deg H_{4,5}\le56$ in independent variables.
\end{corollary}
\begin{proof}
For $S\ne0$, a zero gives the rational root $M-T/S$; by
Theorem~\ref{thm:square} all $A_i$ are integer squares and the root
is an integer. Hence $T/S\in\Z$. Conversely choose a signed sum of
integer roots and add $T/S$ to form $M$. The product representation
has factors $SM-T-S\sum \varepsilon_i p^{i-1}A_i^{(m-1)/2}\sqrt{A_i}$,
which proves the polynomial extension and degree estimate. In
particular at $S=0$ every factor is $-T$.
\end{proof}

Combining Corollary~\ref{cor:H} with Theorem~\ref{thm:div} gives
replacements for both of Sun's packaging gadgets, at the level of
the exact square/divisibility equivalences. If the source square inputs are
polynomials of maximum degree $a$, and its new modulus and numerator
are $\Sigma,\Theta$ of degrees $s,t$, then a fully explicit degree
majorant for the resulting $H_{4,5}$ is
\[
 16\max\{s+1,t,s+5a/2\},\quad
 s\le\deg\xi+5\deg S_0,\qquad
 t\le5\max(\deg T_0,\deg S_0+1).
\]
The last witness $M$ must remain integral; merely allowing a rational
root loses the divisibility information, e.g. all $A_i=1$,
$S=2,T=1$. The substitution uses the same number of source unknowns.
No new numerical total-degree bound for Sun's entire source polynomial
is asserted without expanding its complete Section 7 coefficient DAG.

\section{A guarded encoder from a known separator}\label{sec:guard}
The project's general finite-algebra encoder \cite{encoder} constructs
a separator by spectral difference polynomials of all guarded generators.
The next theorem separates the two tasks. Whenever an inexpensive
separator of geometric characters is already known, a determinant
coefficient formula extracts only the allowed characters, without
multiplying the separator by the character-dependent guard (which can
destroy injectivity). No new ordinary coordinate is introduced.

\begin{theorem}[Guard extraction]\label{thm:guard}
Let $R=\Q[T_1,\ldots,T_b]$, let $\A$ be a finite free commutative
unital $R$-algebra of rank $N\ge1$, and let $g,\beta\in\A$. Suppose
that at every rational specialization $T=t$, the values of $\beta$
on distinct geometric characters are distinct and nonzero. Then one
explicit polynomial $F(T,Y)$ with a single extra rational coordinate
satisfies
\[
 \exists Y\in\Q:F(t,Y)=0
 \iff \exists \chi:\A_t\to\Q:\ \chi(g)\ne0.
\]
If $d_g,d_\beta$ bound the degrees of multiplication-matrix entries
for $g,\beta$, respectively, then
\begin{equation}\label{eq:guarddeg}
 \deg F\le 2\sum_{r=1}^N\max\{r^2 d_g+r(d_\beta+1),\ N d_g\}.
\end{equation}
Multiplication by a fixed integer clears coefficient denominators.
\end{theorem}
\begin{proof}
Let $M_g,M_\beta$ denote the multiplication matrices. Compute
\begin{align*}
 A(X)&=\det(I_N+X M_g)=1+\sum_{r=1}^N a_r X^r,\\
 C(X,Z)&=\det(I_N+X M_g(ZI_N-M_\beta)),\\
 c_r(Z)&=[X^r]C(X,Z)=\sum_{j=0}^r c_{r,j}Z^{r-j}.
\end{align*}
In particular $c_{r,0}=a_r$. Set
\begin{equation}\label{eq:guardpoly}
 H_r(Y)=1+\sum_{j=1}^r c_{r,j} a_r^{j-1}Y^j,
 \qquad
 F=\prod_{r=1}^N\left(H_r(Y)^2+\sum_{s=r+1}^N a_s^2\right).
\end{equation}
At a specialization, all multiplication matrices are simultaneously
triangularizable over $\overline\Q$. The diagonal entries are the
character values with positive multiplicities, whether or not the
algebra is reduced. Let $r_*$ be the number of nonzero guard values
counted with multiplicity. Then $a_{r_*}$ is the product of these
nonzero values (if $r_*=0$ use $a_0=1$), $a_s=0$ for $s>r_*$, and
$c_s(Z)=0$ for $s>r_*$. At $r=r_*>0$,
\[
 c_r(Z)=a_r\prod_{\chi(g)\ne0}(Z-\chi(\beta)),\qquad
 H_r(Y)=\prod_{\chi(g)\ne0}(1-a_rY\chi(\beta)),
\]
where both products count multiplicities. For $r>r_*$, $H_r=1$.
For $r<r_*$ the tail sum has the nonzero square $a_{r_*}^2$.
If $r_*=0$ no factor of $F$ can vanish. Thus a rational zero is
exactly a rational value $1/(a_rY)$ of the separator at an allowed
character; the nonzero hypothesis ensures the converse as well.
All Galois conjugates of such a character have the same separator
value, so separation forces the character itself to be rational.

For the degree bound, determinant minors give
$\deg a_r\le r d_g$ and $\deg c_{r,j}\le r d_g+j d_\beta$, since
each coefficient in $Z$ uses exactly $j$ factors from $M_\beta$.
The corresponding term of $H_r$ has total degree at most
$j(r d_g+d_\beta+1)$. Sum the doubled maxima, including tail squares,
to obtain \eqref{eq:guarddeg}. No extra root, guard reciprocal, or
coefficient-stratum coordinate is needed.
\end{proof}

\begin{remark}
The constant term one in $H_r$ is crucial. Directly using $c_r(Y)$
with the tail vanishing would accept an identically zero $c_r$ at
strata $r>r_*$. The construction above automatically rejects them.
For an arbitrary number field $K$, replace the sum of squares in each
factor by an anisotropic norm form in its $1+N-r$ arguments. For instance
choose a field extension of that degree via an Eisenstein polynomial
at a finite prime and take the norm of the linear combination in a
basis. A norm is zero only for the zero element. This gives the same
one-coordinate equivalence over $K$ with a corresponding degree bound
using that norm-form degree rather than two.
\end{remark}

For independent rational quadratic roots one may use the separator
$\beta=1+\sum_{i=1}^k p^i A_i^{(m-1)/2}\sqrt{A_i}$ with odd $m>k$.
The difference argument of Theorem~\ref{thm:square} proves separation;
the same valuation classes including the constant tag zero show that
$\beta$ never vanishes. Theorem~\ref{thm:guard} therefore encodes any
polynomial guard on these characters, with no generic spectral
separation step. The following section handles the actual two-layer
project architecture, where pole radicands are not independent of
point-root choices.

\section{Two-layer point and pole encoder}\label{sec:receive}
\subsection{Algebra and valuation bounds}
Use the rank-four project point algebra over $\Q[T]$ for each square
abscissa chart. Over $K=\Q(\sqrt2)$ it is a tensor of the two quadratic
root algebras
\[
 j_i^2=A_i+B_i\sqrt2=:a_i,\qquad (j_i')^2=A_i-B_i\sqrt2=:a_i'.
\]
Here the rational basis is $1,U_i,V_i,W_i$, with
$j_i=(U_i+V_i\sqrt2)/2$, $j_i'=(U_i-V_i\sqrt2)/2$,
$W_i=j_i j_i'$. The stated multiplication table is precisely this
invariant algebra, including at zero roots. Let there be $b$ point
blocks, $n=2b$ radical slots, and $N_0=4^b$ point-algebra rank. Let the
terminal layer consist of $c$ monic quadratic roots $q_j^2=C_j$ whose
radicands lie in the point algebra (they may depend on all point roots,
but not on one another). The full algebra has rank $N=N_0 2^c$.
Write $C_j=\sum_\alpha C_{j,\alpha}e_\alpha$ in the rational tensor
point basis. Its coefficients are polynomials in the ordinary base
coordinates, after the receiving source's monic denominator clearing.
All guards, including those denominators, are included in one element
$g$; the proof permits $g$ to vanish on any subset of characters.

For the actual square-abscissa charts, $a_i=h_i^4-(3-\sqrt2)^2$,
where $h_i\in K$ has rational components. It is never zero at a rational
base specialization: if $h_i^4=(3-\sqrt2)^2$ then
$h_i^2=\pm(3-\sqrt2)$, whose norm would imply $\Nm(h_i)^2=7$.
Thus $N_i=A_i^2-2B_i^2$ is always nonzero. This property is used for
the denominator of the scale below; it does not impose nonzero pole
roots. For other first-layer architectures with zero radicands, a
finite branch on which $N_i$ vanish suffices: omit their scale factors,
impose the zero equalities by a sum of squares, include the other $N_i$
in the guard, and take the finite product over branches.

Fix $p=7$, and embed $K$ into $\Q_7$ with $\sqrt2\equiv-3\pmod7$.
Then $\pi=3+\sqrt2$ satisfies $v_7(\pi)=1$, $v_7(\bar\pi)=0$ because
$\pi\bar\pi=7$. All point radicands at a rational base specialization
lie in $\Q_7$. Every point-root value lies in the same compositum of
quadratic extensions of $\Q_7$, with ramification index at most two.
Indeed the squareclass group of an odd local field is generated by
a uniformizer and a nonsquare residue unit: the unit quadratic extension
is unramified, the uniformizer quadratic extension has ramification
index two. Therefore every nonzero value $C_j$ has valuation in
$\tfrac12\Z$, and every nonzero terminal root in $\tfrac14\Z$.

For rational $x$ define
\[
 A_p(x)=\frac1{1+px^2},\qquad
 B_p(x)=\frac{px}{1+px^2}\quad (x\ne0).
\]
Unequal parity of the two denominator valuations gives
\begin{equation}\label{eq:scaleval}
 v_p(A_p(x))\ge\max(0,-v_p(x)),\quad
 v_p(B_p(x))\ge |v_p(x)|.
\end{equation}
The first expression is one at zero; neither denominator can vanish.
Put
\begin{equation}\label{eq:scale}
 W=\prod_{i=1}^b A_p(A_i)A_p(B_i)B_p(N_i)
     \prod_{j=1}^c\prod_\alpha A_p(C_{j,\alpha})=\frac{W_N}{W_D}.
\end{equation}
Use the literal numerator and denominator from this product, without
cancellation. At a rational specialization $W$ is nonzero. Let
$D_*$ be the sum of all the right-hand valuation bounds in
\eqref{eq:scaleval} used in this product (zero contributes zero).
Then $v_p(W)\ge D_*$, and for each nonzero first-layer radicand
$|v_p(a_i)|,|v_p(a_i')|\le D_*$. In fact their valuations are bounded
below by $-\max(0,-v_p(A_i),-v_p(B_i))$ and above by $v_p(N_i)$
minus the lower bound for the conjugate. Every tensor basis element
$e_\alpha$ is a sum of products of at most $n$ first-layer roots
with unit coefficients (at 7, $2$ and $\sqrt2$ are units), so has
valuation at least $-nD_*/2$. It follows that
\begin{equation}\label{eq:terminalbound}
 v_p(C_j)\ge-(n+2)D_*/2,\qquad
 v_p(q_j)\ge-(n+2)D_*/4,
\end{equation}
including zero values with valuation $+\infty$. These bounds hold for
\emph{all} geometric characters, including those excluded by $g$.

\subsection{The separator and its proof}
\begin{theorem}[Two-layer tagged separator]\label{thm:two}
Choose odd integers $m>n$ and $l>c$, put
$\lambda_i=\pi^{2i-1}\bar\pi^{2i}$ and $h_*=m+l(n+2)$, and define
\begin{align*}
 b_1&=1+\sum_{i=1}^b\bigl(\lambda_i j_i^m+\bar\lambda_i(j_i')^m\bigr),\\
 b_2&=\sum_{j=1}^c p^j q_j^l,\qquad E=p^{n+1}W^{h_*},\\
 \beta&=W_D^{h_*} b_1+p^{n+1}W_N^{h_*} b_2.
\end{align*}
The element $\beta$ belongs to the full rational polynomial algebra
and has distinct, nonzero values on all geometric characters at every
rational specialization. With Theorem~\ref{thm:guard} it yields a
one-coordinate guarded finite-algebra encoder for the point and pole
tower, even when terminal roots vanish or the tower is nonreduced.
\end{theorem}
\begin{proof}
The two terms of each point pair are conjugate under the scalar
involution $\sqrt2\mapsto-\sqrt2$, which interchanges $j_i,j_i'$;
thus $b_1$ descends to the rational algebra. The other terms are
rational polynomials in its generators. The factor $W_D^{h_*}$ is
nonzero at all rational specializations, so it suffices to prove
separation for $b_1+Eb_2$.

The point terms have valuations $k+(m/2)v_p(a)$ with tags
$k=1,\ldots,n$ and $v_p(a)\in\Z$. Their valuation classes are
pairwise distinct, also distinct from the constant one (tag zero),
because $m$ is odd and $m>n$. Consequently $b_1$ never vanishes and
$v_p(b_1)\le0$. At two different point characters the difference
$b_1(\chi)-b_1(\psi)$ is a nonempty sum of changed-sign terms.
It has a unique term of least valuation by the same argument (the
factor 2 is a unit), and
\[
 v_p(b_1(\chi)-b_1(\psi))\le n+mD_*/2.
\]
By \eqref{eq:terminalbound}, the valuation of each $b_2$ value and
hence any difference of two such values is at least
$-l(n+2)D_*/4$. Meanwhile
$v_p(E)\ge n+1+[m+l(n+2)]D_*$. Hence the terminal perturbation has
valuation strictly greater than $n+mD_*/2$; it cannot cancel the
point difference or the value of $b_1$ itself. This proves separation
between different point characters and nonvanishing everywhere.

For two characters over the same point character, the terminal
radicands are identical. If the characters are distinct, their
terminal values differ in sign at some nonzero root. The difference
of the $b_2$ values has nonzero summands with valuations
$j+l v_p(q_j)\in j+(l/4)\Z$. Equality between distinct tags would
imply $4(i-j)\equiv0\pmod l$, impossible because $l$ is odd and
$l>c$. The unique least valuation prevents cancellation. Since $E$
is nonzero, this proves separation within each point fiber. All
arguments use only distinct geometric characters, not nilpotent
parts of the algebra. The determinant proof already covers nilpotents.
\end{proof}

\subsection{Explicit degree ledger and project application}
Let $\delta\ge1$ be any integer majorant of the following finite list:
\begin{itemize}
\item total degrees of the $A_i,B_i$ and all polynomial basis coefficients
$C_{j,\alpha}$ after monic denominator clearing;
\item total degrees of all entries of the full multiplication matrices
of $j_i,j_i'$ over $K$ and $q_j$, and of $g$ over $\Q$.
\end{itemize}
All these quantities are computable from the receiving coefficient
DAG and its explicit multiplication tables. This definition of
$\delta$ is not the $g,g\theta$ majorant of the older generic encoder.
No numerical evaluation of $\delta$ for the entire universal source
is being assumed. The literal scale numerator degree is at most
$2b\delta$ and denominator degree at most $(8b+2cN_0)\delta$:
for each point norm $N_i$ the degrees are at most $2\delta$, and its
three denominator factors cost $2\delta+2\delta+4\delta$; each terminal
coefficient contributes at most $2\delta$. Matrix powers give
\begin{equation}\label{eq:beta-degree}
 d_\beta\le B_*\delta,\quad
 B_*=[m+l(n+2)](8b+2cN_0)+\max(m,l).
\end{equation}
Thus the degree of the guarded encoder is at most $K_*\delta$, where
\begin{equation}\label{eq:degreeK}
 K_*=2\left[\frac{N(N+1)(2N+1)}6
 +(B_*+1)\frac{N(N+1)}2+N^2\right].
\end{equation}
This follows by majorizing the maximum in \eqref{eq:guarddeg} by the
sum and using $\delta\ge1$. It is a deliberately conservative bound.

The direct-index transfers \cite{index} have $2k+1$ independent point
blocks and two or three terminal roots. With Sun's ten-integer source,
$k=10$, $b=21$, $n=42$, $m=43$, $N_0=2^{42}$. Use $l=3,5$ respectively.
The exact integer arithmetic in \eqref{eq:beta-degree}--\eqref{eq:degreeK}
gives the following ledger:
\begin{center}\small
\begin{tabular}{lrrrrr}\toprule
Route & $c$ & $N$ & $h_*$ & $B_*$ & $K_*<$\\\midrule
48 unknowns & 2 & $2^{44}$ & 175 & 3078632557802243 & $2^{140}$\\
50 unknowns & 3 & $2^{45}$ & 263 & 6940117394566339 & $2^{143}$\\\bottomrule
\end{tabular}
\end{center}
The single-branch degree bounds are therefore strictly below
$2^{140}\delta$ and $2^{143}\delta$. The original transfer takes the
product of $k+1=11$ common-shift branch polynomials. If each branch
uses the same majorant $\delta$ (take their maximum), the bounds on
the entire instance polynomial are strictly below
$11\cdot2^{140}\delta$ and $11\cdot2^{143}\delta$, respectively.
All coefficients can be obtained by the finite formulas above.
No spectral differences of guarded generators need be formed. This
improves the generic encoder's rank dependence, from a bound of order
$N^7$ for fixed generator count and entry-degree majorant to the
explicit bound \eqref{eq:degreeK} with $B_*=O(N_0)$ for a fixed
number of stages and tags (with the displayed polynomial dependence
on $b,c$). A numerical comparison to every older specialized encoder
requires their actual coefficient DAGs and is not asserted.

The 48/50 variable counts remain exactly as in \cite{index}:
$2(2k+1)$ point parameters, one common slope, $2c$ pole parameters,
and the single encoder. The shorter logarithm in Section~\ref{sec:precision}
changes only the terminal radicands and the guards; the algebra shape
and nonzero point-norm proof remain valid. Thus the two improvements
compose. The fully specified route continues to use its full-family
modular-height premise, not the unrelated fixed $L=143$ construction.

For the 25/27/11 definitions, the rank-four point prefix and independent
terminal quadratic shape are also promising for this method, but their
current v88 specialized positive/projective encoders already have much
smaller degree formulas than the generic spectral encoder. Replacing
those encoders without comparing their actual degrees would not be
justified. The 36-unknown source includes a mixed-integer rank-sixteen
wrapper preceding the terminal pole. The two-layer theorem does not
apply literally to that wrapper. One can split its rank-eight signed
sum into three independent radicals and its final quadratic into a
third layer; the general multi-stage idea described below is a
separate checkpoint, not a claimed 36-route degree record.

\section{Further transfers and the remaining variable-count problem}
\subsection{More layers: a precise extension principle}
The two-layer argument has a general form. Suppose a tower has finitely
many stages, each adjoining several independent square roots of elements
of preceding stages. At a fixed odd prime, the splitting field of
all roots through stage $r$ has ramification index dividing $2^r$:
inductively each new stage is contained in the compositum of quadratic
extensions of the previous local field, whose ramification index is
at most two. Thus differences of two sign choices at the first stage
where characters differ can be tagged modulo $m_r/2^r$ with odd
$m_r$ larger than the number of tags in that stage. To make later
stages negligible, one needs a rational nonzero scale with sufficiently
large valuation relative to both upper and lower valuations of all
nonzero radicands.

Such a scale can be constructed by coefficient strata without computing
spectral differences. For a radicand $C$ in a preceding finite free
algebra of rank $d$, compute the monic characteristic polynomial
$X^d+a_1X^{d-1}+\cdots+a_d$ of multiplication by $C$.
At a specialization let $s$ be its last nonzero coefficient index
(or $s=0$). Put $D^- = \sum_j\max(0,-v_p(a_j))$ with zero contributing
zero. Every eigenvalue has valuation at least $-D^-$: otherwise the
leading term of the monic equation has strictly smaller valuation
than every other term. If $s>0$, the product of the nonzero eigenvalues
counted with multiplicity is $\pm a_s$, so each nonzero eigenvalue
has valuation at most $|v_p(a_s)|+(d-1)D^-$. It follows that
\[
 B_p(a_s)\prod_{j=1}^d A_p(a_j)^d
\]
has valuation at least a common absolute-valuation bound for all
nonzero eigenvalues. For $s=0$ omit the $B_p$ factor. Take the product
of these scales over radicands, branch on the finite coefficient strata
using equalities $a_j=0$ for $j>s$ and guard $a_s\ne0$ for $s>0$.
One can then choose a common $E$ of valuation greater than
$\max_r m_r'$ plus $\max_r m_r$ times the sum of these bounds, where
$m_r'$ counts the tags at stage $r$, and form the hierarchical sum
$1+\sum_r E^{r-1} b_r$ with $b_r$ the tagged sum at that stage.
At the first stage where two characters differ, its difference has a
unique least-valuation summand, and all later differences have larger
valuation after factoring out $E^{r-1}$. At the first stage, the constant
one and nonzero tags give nonvanishing by the same residue calculation;
later terms cannot cancel it. Clear rational scale denominators,
apply Theorem~\ref{thm:guard} in each allowed coefficient stratum, add
the stratum equalities by sums of squares, and take the product over
strata. Invalid strata with a vanishing last coefficient are rejected
by the guard; strata with a nonzero tail are rejected by the equalities.
This proves a finite one-coordinate encoder for such towers, allowing
zero radicands. The number of strata may be large; the sharper bound
of Section~\ref{sec:receive} avoids them for the actual two-layer
architecture. A numerical degree ledger for the mixed wrapper is left
open rather than inferred from the two-layer formula.

\subsection{Integrality absorption and its limits}
There is a useful generic coordinate-saving criterion, independent of
Sun's particular source. Suppose a rational integer-tuple detector
certifies $r$ ordinary inputs integral, and a source further adjoins
monic algebraic equations over the ring generated by those integer
inputs, with the additional roots required rational. Then all these
roots are integral: $\Z$ is integrally closed in $\Q$, or simply apply
the rational-root theorem successively. They require no additional
integer predicates. They can be included in the same joint finite
algebra as the detector, provided all relations, leading coefficients
and guards are encoded with the correct selected-character semantics.
Over a number field the identical assertion holds with $\OK$ and $K$.
For example in the $3r+8$ selective ordinary-$abc$ tuple architecture
\cite{selective}, an applicable triangular monic source can be coded
with $3r+8$ ordinary coordinates rather than testing its additional
monic roots as independent integer inputs. This is a criterion, not
a new universal source with fewer than ten integers.

Two tempting misapplications fail. First, Sun's final homogeneous
merger has leading coefficient $S^{2^k}$ rather than one. Rescaling to
$m'=S M$ makes a monic equation for $m'$, not for $M$. A rational
character does not imply $M\in\Z$ (see the explicit $S=2,T=1$
example following Corollary~\ref{cor:H}). That integer quotient carries
essential divisibility. Second, in the direct-index transfer actual
integer labels are elliptic indices, not ordinary rational base
coordinates of its finite algebra. Eliminating an integer root by
only testing that a monic polynomial has roots modulo the selected
primes would accept the intersective quintic with no global rational
root. Neither argument removes a certified source input.

\subsection{Status and next steps}
\begin{center}\small
\begin{tabular}{p{0.46\linewidth}p{0.45\linewidth}}\toprule
Proved improvement & Scope and remaining restriction\\\midrule
Degree-five intersective merger; simple local roots; all number fields
 & Lower degree at the coding-gadget level. No new integer unknown saved.\\
Tagged simultaneous-square and homogeneous merger
 & $\deg J_4\le40$, $\deg H_4\le56$ in independent inputs; keep integer quotient.\\
Precision $\lfloor D/2\rfloor+1$ and short logarithm
 & Direct 48/50 transfer with same arithmetic premises.\\
One-coordinate guarded tagged encoder, two-layer degree ledger
 & Actual direct-index architecture; no full coefficient expansion.\\
Multi-stage and monic absorption criteria
 & Need explicit new source/degree ledger for further counts.\\\bottomrule
\end{tabular}
\end{center}
The main next step for the quartic frontier is an actual low-quadratic-cost
source with fewer certified integral inputs, not merely a lower-degree
intersective polynomial. The next step for unrestricted degree is to
rebuild the direct-index coefficient DAG with the shorter logarithm and
new encoder, and compare exact numerical degree bounds to the present
specialized positive/projective encoders. For number fields the
intersective polynomial and local interpolation transfer directly, but
neither proves stable rank, a modular height inequality, or global
rational points of the required pole fibers. The new results do not
establish the outstanding 21/9 defining-coordinate targets or the
quartic target near 263. No conjecture is silently promoted to a theorem.

\section*{Verification and provenance}
The accompanying exact Python script checks the quintic expansion,
discriminant and resultant, existence of simple roots at all primes
in a finite test range, exact local lifting samples, the matrix
interpolation permutation at sample primes (including the 2-adic
threshold), the tagged square valuation arithmetic, the guarded
coefficient formula on small algebras with zero roots and excluded
characters, truncation inequalities, and the degree-multiplier ledger.
The executed script returned 303 prime-root checks, 219 exact local
lifts, six number-field fixtures, five scalar and 98 Lucas permutations,
4,000 cutoff checks, 4,800 tag checks, seven guarded tensor examples,
seven distinct characters in a two-layer example, and two degree ledgers.
These are finite regression checks; the proofs in the body establish
the general algebraic statements. The preceding audit of Sun's source
was used, and the repository's receiving proof files were consulted.
During this continuation the arXiv full-text endpoint failed on repeated
retrieval attempts after the abstract/introduction was retrieved;
we therefore do not assert a fresh line-by-line recompilation of the
entire Section 7 source. The precise new gadget substitutions and
transfer proofs above do not depend on such a recompilation. No
repository modifications or full CI runs are claimed in this note.

\begin{thebibliography}{99}
\bibitem{sun26} Z.-W. Sun, \emph{Ten unknowns for Hilbert's tenth problem over the integers}, arXiv:2609.01594v2, 2 September 2026.
\url{https://arxiv.org/abs/2609.01594v2}.
\bibitem{repo} $\Z\Q$ Diophantine Project, \emph{Polynomial presentations of the integers over the rationals}, Jose Brox's improvement repository, README, accessed 24 November 2026.
\url{https://github.com/JoseBrox/diophantine-z-q_improvements}.
\bibitem{index} Project, \emph{Direct polynomial testing on elliptic indices}, v111, \url{https://github.com/JoseBrox/diophantine-z-q_improvements/blob/main/undecidability/direct_index_transfer_v111/README.md}.
\bibitem{v124} Project, \emph{Continuation v124: exact normalized moments}, \url{https://github.com/JoseBrox/diophantine-z-q_improvements/blob/main/docs/results_20261124_v124.md}.
\bibitem{ellheight} Project, \emph{Effective elliptic sums}, v10, \url{https://github.com/JoseBrox/diophantine-z-q_improvements/blob/main/papers/effective_sums/effective_elliptic_sums_20261108_v10.tex}.
\bibitem{encoder} Project, \emph{One rational coordinate for a finite family of algebraic witnesses}, 23 November 2026, \url{https://github.com/JoseBrox/diophantine-z-q_improvements/blob/main/undecidability/papers/joint_finite_algebra_encoder.tex}.
\bibitem{selective} Project, \emph{One selective pole for all integer inputs}, v110, \url{https://github.com/JoseBrox/diophantine-z-q_improvements/blob/main/undecidability/abc_selective_pole_v110/README.md}.
\bibitem{bb} M. Bhargava and A. Burungale, arXiv:2611.16130v2, Lemma 2.2 and Theorem 2.3, \url{https://arxiv.org/abs/2611.16130v2}; application in the project's \emph{Product poles and lower degrees in the few-variable constructions}, v88, \url{https://github.com/JoseBrox/diophantine-z-q_improvements/blob/main/papers/few_variables25/bb_product_degree_transfer_20261123_v88.tex}.
\end{thebibliography}
\end{document}
